\label{sec:reading}
The central obstacle of this note is quantitative. The exact first-order chain reaches $3.5\cdot10^{-8}$ and, with
the Euler--Stein constant, $6.6$--$7.2\cdot10^{-9}$ on every network tried; what it costs is the memory of the third
cumulant, $12n^3$ per source-age and layer, $0.5$--$0.7B$ in all, against a floor of $0.1B$. Section~\ref{sec:measure}
shows that this memory is not compressible in the Frobenius sense (Hadamard obstruction), not sparsifiable in the
hub index (flat importance), and not regenerable from the current layer (new shape at every layer); only the
oldest ages collapse onto a shared basis. The eight papers below were read for what they say, between the lines,
about an obstacle of this kind: a sum of many comparable terms whose readout is diagonal in a hidden index, needed
to a relative precision of $10^{-2}$, under a budget linear in the width of the state.

\subsection{What the papers say}
\emph{Hierarchic flows} \cite{LM}. A high-dimensional non-Gaussian probability is estimated and sampled through a
hierarchy of conditionally Gaussian layers: at each scale the variance is a function of the \emph{modulus} of the
coarser scale, and the long-range non-Gaussian dependence of the field is carried entirely by interactions that
are \emph{local in the hierarchy}: a modulus at one scale coupled to the field at a neighbouring scale. The number
of parameters is polylogarithmic in the dimension because the couplings are between hierarchy levels, not between
coordinates. The robustness of the modulus ($q=1$ rather than the square) is a point the paper makes at length.

\emph{Many Hamiltonians are sparsifiable} \cite{BBP} and \emph{optimal sparsifiers for Minkowski sums and sums of
seminorms} \cite{BBCPRZ}. A sum of $m$ terms can be replaced by a reweighted subset of $O(n/\varepsilon^2)$ terms
that preserves the quadratic form, or the sum of seminorms, uniformly over all inputs in dimension $n$; the proof
is importance sampling under a non-redundancy condition, or, without logarithmic losses, the Reis--Rothvoss
weight-perturbation scheme. Two statements matter here. The size of a sparsifier is governed by the
\emph{dimension of the input}, not by the number of terms, and it carries $\varepsilon^{-2}$. And there are sums
that cannot be sparsified at all: sums of tensor-product (rank-one) terms of comparable weight, whose quadratic
form is identity-like, need all their terms (Theorem~1.5 of \cite{BBP}); random terms, by contrast, sparsify.

\emph{Samplizer} \cite{WZ}. Any quantum algorithm making $Q$ queries to the reflection $I-2|\psi\rangle\langle\psi|$
about an unknown state can be run from $\Theta(Q^2/\delta)$ copies of the state, and this quadratic overhead is
optimal: the reflection, a rank-one perturbation of the identity, is built from samples by density-matrix
exponentiation one small step at a time, each step with its fresh copy. The application estimates the angle
between two states from the spectrum of the \emph{product of two reflections} ($e^{i(\pi\mp2\gamma)}$) by phase
estimation: $\Theta(1/\varepsilon)$ queries, $\Theta(1/\varepsilon^2)$ samples, both optimal. Coherent access to
an object and sample access to it are separated by exactly a square.

\emph{Adaptive phase estimation} \cite{LdW}. Without a promise, adaptive (sequential) and non-adaptive (parallel)
phase estimation need the same number of applications of the unitary; with a promise that the phase lies in a
small known set, sequential strategies save a factor that is at most $2$. The lower bound for non-adaptive
strategies is a Farkas certificate: a nonnegative trigonometric polynomial exhibits the infeasibility of any initial
state. Prior knowledge about the answer buys a constant, and the certificate of what it cannot buy is an explicit
dual object.

\emph{Efficient quantum Hermite transform} \cite{JISBJ}. The harmonic-oscillator evolution factorises through
$\mathfrak{sl}_2$ (the number operator and the two ladder operators), which gives an exact fast-forwarding of the
Gaussian (Ornstein--Uhlenbeck) semigroup: any time is one application, because the semigroup is diagonal in the
Hermite basis with multiplier $e^{-kt}$ on degree $k$. The same structure gives Hermite sampling and a Gaussian
Goldreich--Levin algorithm that finds the heavy Hermite coefficients of a function from queries.

\emph{Learning linear threshold functions under the Gaussian} \cite{KNdW}. The Hermite expansion of a ridge
function is a symmetric tensor power of its direction,
$h_k(\langle w,x\rangle)=\sum_{|\alpha|=k}\sqrt{k!/\alpha!}\,w^\alpha h_\alpha(x)$: the degree-$k$ Hermite
tensor of a ridge function is rank one along $w$. A threshold function has Hermite weight $\Theta(K^{-1/2})$ in
every degree band $[K,2K]$, and its gradient points along $w$, so $O(\log(n/\varepsilon))$ gradient queries
identify it, against polynomially many samples.

\emph{Counting paths with an exterior algebra} \cite{ext}. A sum over exponentially many paths of products of
per-vertex factors is estimated by a dynamic program whose state lives in a $2^k$-dimensional algebra (one
exterior coordinate per subset of $[k]$, one auxiliary row of length $D=k$ each): the identity $g\wedge g=0$
removes every walk with a repeated vertex exactly, independent Gaussian vectors make every surviving term
contribute $k!$ in expectation, and independent Gaussian matrices $N(0,1/D)$ attached to (position, vertex)
pairs give path signatures whose products preserve second moments and whose fourth moment grows only like
$(1+2/D)^{k}$ (Lemma~6.3 there, from Rakhshan--Rabusseau). The matrices must be \emph{shared} across all partial
paths through the same vertex for the program to stay small; the moment bound is proved for exactly that reuse.
Trees are counted by pinning a separator vertex and multiplying the component states. The relative second moment
of one trial is $O(k^2)$, so $O(k^2/\varepsilon^2)$ trials give relative error $\varepsilon$.

\subsection{Four laws they hold for the obstacle}
\paragraph{(L1) The hub sum is in the non-sparsifiable class.} Each birth block is $\sum_j a_j\otimes e_j\otimes c_j$,
one rank-one term per hub neuron $j$ with the hub's own coordinate vector as middle leg (Section~\ref{sec:births}),
and at He initialisation the terms have comparable weight and are orthogonal in the slice inner product (the
least-squares reweighting factor of the kept columns is $1.000$). This is the identity-like sum of
Theorem~1.5 of \cite{BBP}: no reweighted subset reproduces it. The measurement agrees ($q=\tfrac12$: $\times12$;
$q=\tfrac14$: $\times43$), and the sparsifier size $O(n/\varepsilon^2)$ at the precision the chain needs,
$\varepsilon\approx10^{-2}$, exceeds the $2n$ terms the block has. In the language of \cite{stage8} the
transverse measure of the hub foliation is uniform: there is no atom of dominant weight to keep.

\paragraph{(L2) Nothing stochastic reaches the precision.} The memory enters the final mean through slices that must
be right to about $10^{-2}$ of themselves, and the slices are $O(\rho)\approx3\%$ residuals of $O(1)$ objects.
Any estimator that uses randomness to avoid the width of the state (probes for a Laplacian, sketches of the hub
index, the exterior-algebra signatures) pays the sample law $\varepsilon^{-2}$: the samplizer's quadratic gap and
the $O(k^2/\varepsilon^2)$ trials of \cite{ext} are the same statement, and here it means $10^{4}$--$10^{6}$
probes against a budget of $2n$ chain evaluations. This is the capacity obstruction of Section~\ref{sec:capacity}
again, now with its lower bound: coherent (exact) propagation costs the width of the state and no variance; sample
access costs the square of the precision. For a trilinear readout the obstruction is sharper than for a bilinear
one: a Gaussian sketch of the hub index has zero expectation on every trilinear form (odd moments vanish), so even
an unbiased sketch needs third-moment-faithful signs (cube roots of unity) before the variance question arises.

\paragraph{(L3) Reorganising the computation saves a constant at most.} The chain applies $W_l$ to a state of width
$2n$ per age. The Gaussian weights are a promise about that state (its spectrum, its mean-field part), and
\cite{LdW} is the clean model of what a promise buys when the operator applications are the cost: a constant
factor, never the order of growth, with a dual certificate for the rest. The identity-leg accounting and
Strassen ($24n^3\to12n^3\to10.5n^3$ per source-age) are that constant. The dual certificate here is the
orthogonality of the hub terms in (L1).

\paragraph{(L4) The crossover law of the hub-diagonal readout.} A source of age $a$ has legs in the range of the
cocycle $T_{l\leftarrow l-a}$. If that range is held in a basis $Q$ of rank $r$, the transport of a source costs
$O(n^2r)$ per layer instead of $12n^3$, but the slice $\sum_j (QA)_{ij}(QB)_{ij}(QC)_{kj}$ is a Khatri--Rao
contraction whose core costs $\Theta(nr^3)$: the hub-diagonal readout is cubic in the rank of the subspace. Hence
a source is cheaper in its cocycle range only when
\[
 r^3<n^2,\qquad r<n^{2/3}\approx101\quad(n=1024),
\]
and the rank needed is the one at which the cocycle keeps the slice to $10^{-2}$, which by the Hadamard
obstruction is the $99\%$-energy rank, not the stable rank. \texttt{scripts/diag\_cocycle.py} gives, at
layers $8$--$15$, top-$256$ energy fractions $0.80,\,0.92,\,0.97,\,0.99,\,0.997,\,0.998$ at ages
$1,2,3,4,6,8$: the $99\%$-rank is above $256$ for ages $1$--$3$ and below $100$ from age five or six. The crossover
law therefore predicts a dense window of four ages and a subspace for the rest, which is exactly the lossless
configuration found by scanning (window $4$, $k=384$: $3.69\cdot10^{-8}$) and not the cheaper ones (window $2$:
$\times2.2$). Its cost is $4\times12n^3$ per layer before Strassen against a floor of $12.8n^3$: an exact first-order
chain cannot be brought below $3$--$4$ times the floor by any representation of its memory, because for the ages
that carry the memory the dense Hadamard readout is already the cheapest one.

\subsection{Three readings that point forward}
\paragraph{(P1) Fast-forwarding the Eldan path: the merge constant is a Hermite-degree average.} Eldan's path is an
Ornstein--Uhlenbeck semigroup in the tilt variable, and \cite{JISBJ} is the reminder that such a semigroup is
diagonal in Hermite degree: a defect whose kink content sits at degree $d$ decays along the path like
$(1+\tau)^{-d/2}$ (one kink: the density $(1+\tau)^{-1/2}$ of Section~\ref{sec:capacity}; two kinks: their
product). Against the weight $\tfrac12(1+\tau)^{-3/2}$ this gives, with the normalisation of
Corollary~\ref{cor:eldan},
\[
 a(d)=\frac{1}{d+1}:\qquad a(1)=\tfrac12\ \text{(Gaussian closure: measured $0.535$, $0.500$)},\quad
 a(2)=\tfrac13\ \text{(first-order chain: measured $0.255$, $0.257$)},
\]
the second reduced by the rotation of the defect direction along the path, as on the small network. The constant
is therefore a property of the \emph{order} of the chain, not of the network, which is what the cross-over of
$a_0$ and $a_1$ without loss on both networks shows. A system can fix it offline from the chain's order and never
fit it.

\paragraph{(P2) The memory is a sum over hub-rooted stars, and one of its three legs collapses.} With
$S_{l'}\approx T_{l'\leftarrow0}T_{l'\leftarrow0}^{\!\top}$ the star block born at $l'$ and read at $l$ is
\[
 \sum_{j}\big(T_{l\leftarrow l'}\Phi S_{l'}\big)_{ij}\,\big(T_{l\leftarrow l'}\big)_{ij}\,\big(T_{l\leftarrow l'}w_2S_{l'}\Phi\big)_{kj}
 =\sum_j\sum_{m,m'}T_{l\leftarrow l'}(i,j)\,T_{l'\leftarrow0}(j,m)\,T_{l'\leftarrow0}(j,m')\,A^{(l')}_{im}B^{(l')}_{km'},
\]
a sum over hub vertices $j$ at every layer $l'$ of two paths from the input to the hub and one path from the hub
to the output, the tree-counting situation of \cite{ext} with the hub as separator. The first leg accumulates:
$U_l=\sum_{l'\le l}T_{l\leftarrow l'}D(\Phi_{l'})T_{l'\leftarrow0}$ obeys $U_{l+1}=W_{l+1}D(g_{l+1})U_l+D(\Phi_{l+1})T_{l+1\leftarrow0}$,
one $n\times n$ matrix for all sources at $n^3$ per layer. What does not collapse is the pairing of the hub leg
$T_{l\leftarrow l'}(i,j)$ with the third leg at the same $(l',j)$: the Hadamard pairing is diagonal in the hub and
the hub's age is a label the sum cannot forget. This locates the memory precisely: not in the transported
content (which accumulates) but in the hub-diagonal pairing of the identity leg with the readout leg of the same
birth. The coherent part of that pairing is the mean of $T\circ T$ over the weights, a variance profile
computable at $O(n^2)$, and the rank-one experiments of Section~\ref{sec:measure} measure it to carry $80\%$ of
the memory; the rest is the realisation-specific fluctuation of a Gaussian matrix product, of the relative size
the fourth-moment bound of \cite{ext} predicts for products of $a$ matrices of width $n$, a few percent per age.

\paragraph{(P3) One object: memory, tilt response, defect.} Theorem~\ref{thm:gram} writes the Euler--Stein defect as
$-\langle\nabla^2G,\Gamma(y)\rangle$ with $\Gamma$ built from the first-order tilt responses, and the second-order
response of the Gaussian chain to the $n$ coordinate tilts has, at every layer, the birth structure
$\sum_a (\text{cocycle column})\otimes e_a\otimes(\text{row of }X_l)$: an identity leg born at each layer and
transported by the same cocycle. The third cumulant's memory, the tilt-response tensor and the defect are the
same cocycle-transported identity-leg object read through the same hub-diagonal pairing. Two consequences. The
symbolic Euler--Stein correction (the open task of Section~\ref{sec:measure}) costs what the memory costs, so it
is not an add-on to the exact chain but a second use of the same transported legs, to be built on them. And the
cheap part of each is the same cheap part: the mean-field pairing at $O(n^2)$.

\subsection{What this says the system is}
An exact first-order chain cannot reach the floor: (L4) puts it at three to four times $0.1B$ after every
algebraic saving, and (L1)--(L3) close the other exits. A floor-level estimator with the accuracy of the leaders is
therefore not an exact chain with a compressed memory; it is a chain with a \emph{short} memory, corrected by
derived quantities that cost $O(n^2)$ per layer: the mean-field hub pairing for every age beyond the window (P2),
the radial scalar of Lemma~\ref{lem:scale} (one number, worth a factor $16$), and the Euler--Stein merge with the
constant fixed by the order of the chain (P1), computed from the same legs as the memory (P3). The one measurement
that decides whether this reaches $10^{-8}$ at the floor is whether the defect of a \emph{short-memory} chain is
as explainable by its own coordinate Laplacian as that of the Gaussian closure ($89$--$93\%$) and of the exact chain
($81$--$84\%$): if it is, a window-one chain at the floor with the mean-field tail and the merge lands near
$10^{-8}$; if it is not, the memory is the price and the score is bounded below by the crossover law. That
measurement is $2049$ runs of the window-one chain, which the Modal batch of Section~\ref{sec:measure} does in an
hour; it is the next thing to run.
